% !TeX root = ../main.tex
\section{Access Pattern and Granularity}\label{sec:pattern}
\rubric{Pattern, granularity and read/write sweeps}{35}{Complete, well-controlled bandwidth and latency
comparisons (with \cref{sec:rw-mix}).}

\begin{hypothesis}
  Sequential access should let the prefetcher amortise DRAM latency, giving high bandwidth that
  \emph{falls} as the stride grows and each fetched \qty{64}{\byte} line delivers fewer useful bytes.
  Random access should be near-flat and low across granularities, since every access is a fresh cache
  miss regardless of stride.
\end{hypothesis}

\begin{knobs}
  Fixed       & 1 thread pinned to CPU~2; \qty{256}{\mebi\byte} DRAM-resident region; 100\% reads \\
  Varied      & pattern (sequential, random); byte stride (8, 64, 256, 1024) \\
  Script      & \code{Cache_and_Memory_Performance_Profiling/scripts/run.sh granularity} and \code{latency} \\
\end{knobs}

\begin{figure}[tbp]
  \centering
  \begin{subfigure}[t]{0.49\linewidth}
    \IfDataExists{granularity}{%
    \begin{tikzpicture}
      \begin{axis}[acs, xmode=log, log basis x=2, ymode=log, height=0.85\linewidth,
                   xlabel={Byte stride}, ylabel={Bandwidth (GiB/s)}, legend pos=south west]
        \addplot+[iqr bars] table[x=strideb, y=gib_s_median, iqr of=gib_s, only if={pattern}{sequential}] {data/granularity.csv};
        \addlegendentry{sequential}
        \addplot+[iqr bars] table[x=strideb, y=gib_s_median, iqr of=gib_s, only if={pattern}{random}] {data/granularity.csv};
        \addlegendentry{random}
      \end{axis}
    \end{tikzpicture}}
    \caption{Delivered read bandwidth.}
  \end{subfigure}\hfill
  \begin{subfigure}[t]{0.49\linewidth}
    \IfDataExists{latency}{%
    \begin{tikzpicture}
      \begin{axis}[acs, xmode=log, log basis x=2, ymode=log, height=0.85\linewidth,
                   xlabel={Working set (bytes)}, ylabel={Latency (ns/access)}, legend pos=north west]
        \addplot+[iqr bars] table[x=bytes, y=ns_per_access_median, iqr of=ns_per_access, only if={pattern}{random}] {data/latency.csv};
        \addlegendentry{random}
        \addplot+[iqr bars] table[x=bytes, y=ns_per_access_median, iqr of=ns_per_access, only if={pattern}{sequential}] {data/latency.csv};
        \addlegendentry{sequential}
      \end{axis}
    \end{tikzpicture}}
    \caption{Pointer-chase latency vs.\ size.}
  \end{subfigure}%
  \caption{Access pattern and granularity. Left: sequential read bandwidth falls steeply with stride as
  each \qty{64}{\byte} line yields fewer useful bytes, while random bandwidth is flat and low. Right: the
  prefetcher keeps sequential latency near \qty{5}{\nano\second} even in DRAM, whereas random latency
  tracks the full hierarchy.}\label{fig:pattern}
\end{figure}

\subsection{Analysis}
Sequential bandwidth drops from \qty{14.8}{\gibi\byte\per\second} at unit (8-byte) stride to
\qty{2.4}{\gibi\byte\per\second} at \qty{64}{\byte}, \qty{1.09}{\gibi\byte\per\second} at
\qty{256}{\byte} and \qty{0.93}{\gibi\byte\per\second} at \qty{1024}{\byte}. The knee is at the
\qty{64}{\byte} cache line: once the stride reaches a line, every access still pulls a full line but uses
only one element of it, so \emph{useful} bandwidth falls roughly as $1/\text{stride}$ while the prefetcher
can no longer help. Random access is flat at \qtyrange{0.85}{0.93}{\gibi\byte\per\second} across all
strides---each access is an independent miss, so granularity barely matters. The latency panel is the
same story from the other side: sequential access to a DRAM-resident region costs \qty{5.3}{\nano\second}
(prefetched) while random access costs \qty{84}{\nano\second}, a \num{16}$\times$ gap that is entirely
the prefetcher. This is the mechanism behind the cache-miss results in \cref{sec:cache-miss}.
